\documentclass[11pt,a4paper]{ltjsarticle}
\usepackage{amsmath,amssymb}
\usepackage{graphicx}
\usepackage{booktabs}
\usepackage{array}
\usepackage{luatexja-fontspec}
\setmainjfont{HaranoAjiMincho-Regular}
\setsansjfont{HaranoAjiGothic-Medium}
\linespread{1.05}\selectfont

% Recommended PDF filename:
% アルゴリズム設計論_repo2_2611193_中井平蔵.pdf
% Example build command:
% latexmk -lualatex -jobname=アルゴリズム設計論_repo2_2611193_中井平蔵 main.tex

\begin{document}

\begin{center}
{\LARGE \textbf{アルゴリズム設計論 第2回レポート}}
\end{center}
\vspace{1em}

\begin{flushright}
提出日: \today \\
学籍番号: 2611193 \\
氏名: 中井 平蔵
\end{flushright}

\section*{1.}
Shorのアルゴリズムは、RSA暗号における「大きな合成数の素因数分解は計算量的に困難である」という仮定を破る。
RSA暗号では、公開鍵に含まれる大きな合成数 $N$ が、2つの大きな素数 $p$ と $q$ の積 $N = pq$ であることを利用している。
RSAの安全性は、$N$ が公開されていても、そこから $p$ と $q$ を求める素因数分解が現実的な時間では困難である、
という計算量的仮定に依存している。

しかし、Shorのアルゴリズムを用いると、量子コンピュータ上では素因数分解を多項式時間で解くことができる。
具体的には、$N$ のビット長を $n = \log N$ とすると、
Shorのアルゴリズムは $O(n^3)$ の時間計算量で素因数分解を行えるとされる。
これは、現在の古典コンピュータの素因数分解アルゴリズムよりも大幅に高速である。

\begin{center}
\fbox{\begin{tabular}{@{}l@{}}
したがって、Shorのアルゴリズムは、RSA暗号の安全性の根拠である\\
「素因数分解は計算量的に困難である」という計算量的仮定を破るアルゴリズムである。
\end{tabular}}
\end{center}

\section*{2.}
\subsection*{(1) どのような測定結果がどのような確率で得られるか}
観測量 $Y$ の測定結果は、その固有値として得られる。まず、$Y$ の固有値を求めるために固有方程式
\[
\det(Y - \lambda I) = 0
\]

を解く。
\[
\det \begin{pmatrix}
-\lambda & -i \\
i & -\lambda
\end{pmatrix}
= \lambda^2 - 1 = 0
\]

よって、固有値は
\[
\lambda = \pm 1
\]

である。したがって、測定結果として得られる値は $+1$ または $-1$ の2通りである。

次に、それぞれの固有ベクトルを求める。
\[
\lvert y_+\rangle = \frac{1}{\sqrt{2}}
\begin{pmatrix}
1 \\
i
\end{pmatrix}
= \frac{1}{\sqrt{2}}\lvert 0\rangle + \frac{i}{\sqrt{2}}\lvert 1\rangle
\]

\[
\lvert y_-\rangle = \frac{1}{\sqrt{2}}
\begin{pmatrix}
1 \\
-i
\end{pmatrix}
= \frac{1}{\sqrt{2}}\lvert 0\rangle - \frac{i}{\sqrt{2}}\lvert 1\rangle
\]

状態
\[
\lvert \psi\rangle = \frac{1}{\sqrt{3}}\lvert 0\rangle + \frac{\sqrt{2}}{\sqrt{3}}\lvert 1\rangle
\]

に対して、各測定結果の確率は
\[
p = \left|\langle \text{固有状態} \vert \psi\rangle\right|^2
\]

で与えられる。

初めに、$+1$ が得られる確率は
\begin{align*}
\langle y_+ \vert \psi\rangle
&= \left(\frac{1}{\sqrt{2}}\langle 0\rvert - \frac{i}{\sqrt{2}}\langle 1\rvert\right)
   \left(\frac{1}{\sqrt{3}}\lvert 0\rangle + \frac{\sqrt{2}}{\sqrt{3}}\lvert 1\rangle\right) \\
&= \frac{1}{\sqrt{2}}\cdot\frac{1}{\sqrt{3}}\langle 0 \vert 0\rangle
 + \frac{1}{\sqrt{2}}\cdot\frac{\sqrt{2}}{\sqrt{3}}\langle 0 \vert 1\rangle
 + \left(-\frac{i}{\sqrt{2}}\right)\cdot\frac{1}{\sqrt{3}}\langle 1 \vert 0\rangle
 + \left(-\frac{i}{\sqrt{2}}\right)\cdot\frac{\sqrt{2}}{\sqrt{3}}\langle 1 \vert 1\rangle \\
&= \frac{1}{\sqrt{6}} - \frac{\sqrt{2}i}{\sqrt{6}}
\end{align*}

ここで、$\langle 0 \vert 1\rangle = \langle 1 \vert 0\rangle = 0$ と
$\langle 0 \vert 0\rangle = \langle 1 \vert 1\rangle = 1$ を用いると、
\[
p_+ = \left|\langle y_+ \vert \psi\rangle\right|^2
= \left|\frac{1}{\sqrt{6}} - \frac{\sqrt{2}i}{\sqrt{6}}\right|^2
= \left(\frac{1}{\sqrt{6}}\right)^2 + \left(\frac{\sqrt{2}}{\sqrt{6}}\right)^2
= \frac{1}{6} + \frac{2}{6}
= \frac{1}{2}
\]

である。

同様に、$-1$ が得られる確率は
\begin{align*}
\langle y_- \vert \psi\rangle
&= \left(\frac{1}{\sqrt{2}}\langle 0\rvert + \frac{i}{\sqrt{2}}\langle 1\rvert\right)
   \left(\frac{1}{\sqrt{3}}\lvert 0\rangle + \frac{\sqrt{2}}{\sqrt{3}}\lvert 1\rangle\right) \\
&= \frac{1}{\sqrt{2}}\cdot\frac{1}{\sqrt{3}}\langle 0 \vert 0\rangle
 + \frac{1}{\sqrt{2}}\cdot\frac{\sqrt{2}}{\sqrt{3}}\langle 0 \vert 1\rangle
 + \frac{i}{\sqrt{2}}\cdot\frac{1}{\sqrt{3}}\langle 1 \vert 0\rangle
 + \frac{i}{\sqrt{2}}\cdot\frac{\sqrt{2}}{\sqrt{3}}\langle 1 \vert 1\rangle \\
&= \frac{1}{\sqrt{6}} + \frac{\sqrt{2}i}{\sqrt{6}}
\end{align*}

よって、
\[
p_- = \left|\langle y_- \vert \psi\rangle\right|^2
= \left|\frac{1}{\sqrt{6}} + \frac{\sqrt{2}i}{\sqrt{6}}\right|^2
= \left(\frac{1}{\sqrt{6}}\right)^2 + \left(\frac{\sqrt{2}}{\sqrt{6}}\right)^2
= \frac{1}{6} + \frac{2}{6}
= \frac{1}{2}
\]

となる。したがって、

\begin{center}
\fbox{測定結果は $+1$ と $-1$ であり、それぞれ $\frac{1}{2}$ の確率で得られる。}
\end{center}

\subsection*{(2) 結果の期待値}
期待値 $\langle Y \rangle$ は、測定結果とその確率の積の和として
\[
\langle Y \rangle = \sum_i p_i \lambda_i
\]

で計算できる。よって、
\[
\langle Y \rangle
= \left(\frac{1}{2}\right)(+1) + \left(\frac{1}{2}\right)(-1)
= 0
\]

となる。したがって、

\begin{center}
\fbox{結果の期待値は $\langle Y \rangle = 0$ である。}
\end{center}

\section*{3.}
未知の量子状態を複製する以下の変換
\[
U(\lvert \psi \rangle \lvert 0 \rangle) = \lvert \psi \rangle \lvert \psi \rangle
\]

が存在しないことを示す。

$U$ が存在すると仮定すると、
\[
U(\lvert 0\rangle \lvert 0\rangle) = \lvert 0\rangle \lvert 0\rangle,\qquad
U(\lvert 1\rangle \lvert 0\rangle) = \lvert 1\rangle \lvert 1\rangle
\]

に変換される。ここで、線形性より$U(a\lvert 0\rangle\lvert 0\rangle + b\lvert 1\rangle\lvert 0\rangle)$は以下のように変換される。
\[
U(a\lvert 0\rangle\lvert 0\rangle + b\lvert 1\rangle\lvert 0\rangle)
= a\,U(\lvert 0\rangle\lvert 0\rangle) + b\,U(\lvert 1\rangle\lvert 0\rangle)
= a\lvert 0\rangle\lvert 0\rangle + b\lvert 1\rangle\lvert 1\rangle
\]

また、
\[
\lvert \psi \rangle = a\lvert 0\rangle + b\lvert 1\rangle
\]

としたとき、
\[
U(\lvert \psi\rangle\lvert 0\rangle)
= U((a\lvert 0\rangle + b\lvert 1\rangle)\lvert 0\rangle)
= U(a\lvert 0\rangle\lvert 0\rangle + b\lvert 1\rangle\lvert 0\rangle)
= a\lvert 0\rangle\lvert 0\rangle + b\lvert 1\rangle\lvert 1\rangle
\]

となる。一方、複製できるなら
\[
U(\lvert \psi\rangle \lvert 0\rangle)=\lvert \psi\rangle\lvert \psi\rangle
\]

でなければならないが、
\[
\lvert \psi\rangle\lvert \psi\rangle
=(a\lvert 0\rangle+b\lvert 1\rangle)(a\lvert 0\rangle+b\lvert 1\rangle)
= a^2\lvert 0\rangle\lvert 0\rangle + ab\lvert 0\rangle\lvert 1\rangle + ba\lvert 1\rangle\lvert 0\rangle + b^2\lvert 1\rangle\lvert 1\rangle
\]

となる。よって、
任意の未知の状態を複製する変換 $U(\lvert \psi\rangle\lvert 0\rangle) = \lvert \psi\rangle\lvert \psi\rangle$ を仮定すると、
重ね合わせ状態を入力した際に出力が線形性を満たさなくなり、矛盾する。したがって、

\begin{center}
\fbox{任意の未知の状態を複製する変換 $U(\lvert \psi\rangle\lvert 0\rangle) = \lvert \psi\rangle\lvert \psi\rangle$ は存在しない。}
\end{center}

\end{document}
